% chapters/11/sub-cpmk.tex
\begin{subcpmkbox}
Bab ini menuntaskan \SUBCPMK{2.4} sekaligus menutup \CPMK{2}. Bab ini adalah
bab bertingkat tertinggi pada blok Python, yaitu ranah C5, sehingga penekanannya
bukan pada menulis kode baru, melainkan pada \textbf{menilai} kode yang sudah ada.

\vspace{0.5em}
\textbf{\SUBCPMK{2.4} Mengevaluasi perilaku NPC hasil kombinasi FSM dan A*
terhadap skenario permainan} --- ranah C5, \CPMK{2}, \CPL{03} dan \CPL{04}.

\vspace{0.5em}
{\setlength{\tabcolsep}{4pt}\footnotesize%
\begin{tabularx}{\textwidth}{@{}C{1.9cm} L{4.0cm} L{2.5cm} Y@{}}
\toprule
\headrow \textbf{Indikator} & \textbf{Rumusan} & \textbf{Dibahas pada} & \textbf{Bukti} \\
\midrule
\ind{a} & Menyusun skenario uji dan mencatat hasil pengamatan &
Bagian~\ref{sec:skenario-perbaikan} &
Tiga skenario uji (ruang terbuka, labirin, rintangan dinamis) beserta catatan
pengamatan pada laporan \istilah{playtest}. \\
\addlinespace
\ind{b} & Menilai kualitas perilaku NPC (kelancaran, daya tanggap, beban
komputasi) &
Bagian~\ref{sec:menilai-mutu} &
Tabel penilaian pada beberapa selang hitung ulang, beserta penafsiran ketiga
ukuran mutu. \\
\addlinespace
\ind{c} & Merumuskan perbaikan berdasarkan hasil evaluasi &
Bagian~\ref{sec:skenario-perbaikan} &
Daftar perbaikan yang diusulkan beserta alasan dan cara mengukurnya, dan
perbaikan yang benar-benar diterapkan pada mini proyek. \\
\bottomrule
\end{tabularx}}
\end{subcpmkbox}
